% !TeX encoding = UTF-8
\documentclass[variant=guia,cover=institutional,mode=print]{ua-engineering-v6-article}
% !TeX encoding = UTF-8
\UASetUniversity{Universidad Austral}
\UASetFaculty{Facultad de Ingeniería}
\UASetDepartment{Departamento de Computación}
\UASetCampus{Pilar}
\UASetCareer{Ingeniería Informática}
\UASetCommission{Comisión A}
\UASetProfessor{Equipo docente}
\UASetSemester{Segundo cuatrimestre 2026}
\UASetCourse{Sistemas Distribuidos}
\UASetDate{2026}

\UASetSemester{Segundo cuatrimestre 2026}
\UASetDate{2026}
\UASetClassNumber{07}
\UASetClassTitle{Logs distribuidos y procesamiento de streams}
\title{Clase 07 --- Ficha de aula: historia, progreso y efectos}
\author{\UAProfessor}
\date{\UADate}
\begin{document}
\maketitle
\section{Cómo usar esta ficha}
Complete durante la clase. No busque definiciones perfectas al inicio: escriba una intuición, revísela después del timeline y cierre con una formulación técnica.

\section{1. Intuiciones que luego formalizaremos}
\begin{tabular}{p{0.18\textwidth}p{0.34\textwidth}p{0.34\textwidth}}
\toprule
Concepto & Mi intuición inicial & Definición corregida al final \\
\midrule
Log & \rule{0pt}{1.1cm} & \\
Partición & \rule{0pt}{1.1cm} & \\
Offset & \rule{0pt}{1.1cm} & \\
Consumer group & \rule{0pt}{1.1cm} & \\
Exactly-once & \rule{0pt}{1.1cm} & \\
Event time & \rule{0pt}{1.1cm} & \\
\bottomrule
\end{tabular}


\newpage
\section{2. Identidades y vocabulario del caso CampusShop}

CampusShop confirma una compra y produce un evento. Complete sin mezclar los tres identificadores:

\begin{center}
\begin{tabular}{p{0.26\textwidth}p{0.28\textwidth}p{0.32\textwidth}}
\toprule
Elemento & Ejemplo & Qué identifica \\
\midrule
Operación lógica & \texttt{9F2A} & \rule{0pt}{1.05cm} \\
Evento & \texttt{evt-9F2A-1} & \rule{0pt}{1.05cm} \\
Posición del log & P2, offset 418 & \rule{0pt}{1.05cm} \\
\bottomrule
\end{tabular}
\end{center}

Defina brevemente:

\begin{description}
\item[Broker:] \hrulefill
\item[Topic:] \hrulefill
\item[Partición:] \hrulefill
\item[Offset:] \hrulefill
\item[Committed offset:] \hrulefill
\item[Consumer group:] \hrulefill
\end{description}

\begin{uawarn}
Pregunta de control: ¿dos registros ubicados en P0, offset 42 y P1, offset 42 representan necesariamente el mismo evento? Justifique.
\end{uawarn}

\newpage
\section{3. Particiones y keys}

Eventos: \texttt{OrderCreated}, \texttt{PaymentAuthorized}, \texttt{OrderCancelled} para pedidos 9F2A y 1C7B.

\begin{center}
\begin{tikzpicture}[x=1cm,y=1cm,>=Latex]
\foreach \y/\p in {2/P0,1/P1,0/P2}{
  \node[left,font=\bfseries] at (0,\y){\p};
  \draw[UAIngRule,thick] (0.4,\y)--(12.5,\y);
}
\end{tikzpicture}
\end{center}

\begin{enumerate}
\item Key elegida: \hrulefill
\item Orden que preserva: \hrulefill
\item Orden que no preserva: \hrulefill
\item Hot key posible: \hrulefill
\item Métrica para detectarla: \hrulefill
\end{enumerate}


\newpage
\section{4. Publicación, réplica y ACK}

Suponga una partición con líder A y seguidores B y C. Dibuje el recorrido de un registro y marque cuándo se respondería al productor en cada política.

\begin{center}
\begin{tikzpicture}[node distance=1.1cm,>=Latex,every node/.style={align=center}]
\node[draw=UAIngOrange,rounded corners,thick,text width=2.2cm] (p) {Productor};
\node[draw=UAIngNavy,rounded corners,thick,right=of p,text width=2.2cm] (a) {Líder A};
\node[draw=UAIngOrange,rounded corners,thick,above right=0.35cm and 1.3cm of a,text width=2.2cm] (b) {Seguidor B};
\node[draw=UAIngOrange,rounded corners,thick,below right=0.35cm and 1.3cm of a,text width=2.2cm] (c) {Seguidor C};
\draw[->,thick] (p)--node[above]{publicar}(a);
\draw[->,thick] (a)--(b);
\draw[->,thick] (a)--(c);
\end{tikzpicture}
\end{center}

\begin{tabular}{p{0.24\textwidth}p{0.31\textwidth}p{0.31\textwidth}}
\toprule
Política & Condición de ACK & Falla todavía posible \\
\midrule
Sin esperar & \rule{0pt}{1.0cm} & \\
Solo líder & \rule{0pt}{1.0cm} & \\
Réplicas requeridas & \rule{0pt}{1.0cm} & \\
\bottomrule
\end{tabular}

\vspace{0.35cm}
\textbf{ISR significa:} \hrulefill\\[0.45cm]
\textbf{Trade-off principal:} \hrulefill

\newpage
\section{5. Timeline de offset y crash}

Complete dos veces: primero commit antes del efecto; luego efecto antes del commit.

\begin{center}
\begin{tikzpicture}[x=1cm,y=1cm,>=Latex]
\foreach \y/\lab in {3/Broker,2/Consumer,1/Efecto,0/Offset}{
  \node[left,font=\bfseries] at (0,\y){\lab};
  \draw[UAIngRule,thick] (0.4,\y)--(12.5,\y);
}
\draw[->,thick] (0.4,-0.7)--(12.8,-0.7) node[right]{tiempo};
\end{tikzpicture}
\end{center}

\textbf{Política A:} ¿qué se pierde? \hrulefill\\[0.6cm]
\textbf{Política B:} ¿qué se duplica? \hrulefill\\[0.6cm]
\textbf{Dato durable para idempotencia:} \hrulefill

\newpage
\section{6. Frontera de exactly-once}
Marque con un rectángulo qué componentes participan en su garantía.

\begin{center}
\begin{tikzpicture}[node distance=1cm,>=Latex,every node/.style={align=center}]
\node[draw=UAIngOrange,rounded corners,thick,text width=2.5cm] (in) {Topic input};
\node[draw=UAIngNavy,rounded corners,thick,right=of in,text width=2.5cm] (state) {State store};
\node[draw=UAIngOrange,rounded corners,thick,right=of state,text width=2.5cm] (out) {Topic output};
\node[draw=UAIngNavy,rounded corners,thick,right=of out,text width=2.5cm] (mail) {Proveedor email};
\draw[->,thick] (in)--(state)--(out)--(mail);
\end{tikzpicture}
\end{center}

Complete:
\begin{itemize}
\item Objeto garantizado: \hrulefill
\item Falla cubierta: \hrulefill
\item Efecto fuera de frontera: \hrulefill
\item Mecanismo adicional: \hrulefill
\end{itemize}

\newpage
\section{7. Dual write y outbox}
Dibuje las dos ventanas de falla y luego la solución outbox.

\begin{center}
\begin{tikzpicture}[x=1cm,y=1cm,>=Latex]
\node[draw=UAIngOrange,rounded corners,thick] at (1.5,2){Servicio};
\node[draw=UAIngNavy,rounded corners,thick] at (6,2){Base};
\node[draw=UAIngOrange,rounded corners,thick] at (10.5,2){Broker};
\draw[UAIngRule,thick] (0,0.8)--(12,0.8);
\node[left] at (0,0.8){Orden 1};
\draw[UAIngRule,thick] (0,-0.2)--(12,-0.2);
\node[left] at (0,-0.2){Orden 2};
\end{tikzpicture}
\end{center}

Campos mínimos de outbox: \hrulefill\\[0.5cm]
¿Por qué el relay puede duplicar? \hrulefill\\[0.5cm]
¿Cómo deduplica el consumidor? \hrulefill

\newpage
\section{8. Event time, watermark y eventos tardíos}
Eventos:
\begin{center}
\begin{tabular}{cccc}
\toprule
Evento & Event time & Arrival time & Valor \\
\midrule
E1 & 10:00:40 & 10:00:43 & 10 \\
E2 & 10:02:15 & 10:02:16 & 20 \\
E3 & 10:04:50 & 10:07:10 & 30 \\
\bottomrule
\end{tabular}
\end{center}

Para ventanas de cinco minutos:
\begin{enumerate}
\item Resultado por processing time: \hrulefill
\item Resultado por event time: \hrulefill
\item Con watermark de 2 min, tratamiento de E3: \hrulefill
\item Para fraude, ¿esperaría más o menos? ¿Por qué? \vspace{1cm}
\end{enumerate}


\newpage
\section{9. Diagnóstico de lag y operación}

Complete la interpretación del panel:

\begin{center}
\begin{tabular}{p{0.17\textwidth}p{0.18\textwidth}p{0.22\textwidth}p{0.30\textwidth}}
\toprule
Partición & Ingreso/s & Consumo/s & Hipótesis \\
\midrule
P0 & 10.000 & 10.300 & \rule{0pt}{0.9cm} \\
P1 & 10.500 & 10.600 & \rule{0pt}{0.9cm} \\
P2 & 38.000 & 16.000 & \rule{0pt}{0.9cm} \\
P3 & 9.800 & 10.100 & \rule{0pt}{0.9cm} \\
\bottomrule
\end{tabular}
\end{center}

\begin{enumerate}
\item ¿Qué partición acumula trabajo? \hrulefill
\item ¿Agregar consumidores divide automáticamente P2? \hrulefill
\item ¿Qué diferencia hay entre lag en offsets y edad del evento? \vspace{0.7cm}
\item Proponga tres métricas adicionales. \vspace{0.8cm}
\item ¿Qué podría significar una DLQ creciente? \vspace{0.8cm}
\end{enumerate}

\newpage

\begin{uainfo}
\textbf{Precisión para la ficha:} el identificador del evento y la marca de deduplicación deben persistirse atómicamente con el efecto local. Registrar primero el efecto y después el ID deja una ventana de duplicación.
\end{uainfo}

\section{10. Diseño y exit ticket}
Caso: ranking de Asteria Online.

\begin{tabular}{p{0.27\textwidth}p{0.60\textwidth}}
\toprule
Decisión & Propuesta \\
\midrule
Hecho e ID & \rule{0pt}{0.9cm} \\
Key y orden & \rule{0pt}{0.9cm} \\
Semántica & \rule{0pt}{0.9cm} \\
Late data & \rule{0pt}{0.9cm} \\
SLO & \rule{0pt}{0.9cm} \\
Prueba de falla & \rule{0pt}{0.9cm} \\
\bottomrule
\end{tabular}

\subsection*{Exit ticket}
\begin{enumerate}
\item Diferencia entre event ID y offset.
\item Timeline mínimo que produce duplicación.
\item Frontera concreta de exactly-once.
\item Métrica para detectar una hot partition.
\end{enumerate}
\end{document}
